% !TEX root = ../../../main.tex
\section{Linear functions and maps}
\label{sec:linear-algebra-i:linear-functions-and-maps}

% Sources: supplied lecture photos, quiz-prep p. 5, September 17 slides
% 18--25, and September 24 slides 18--19.
\subsection{Linearity and linear functionals}

\begin{definition}
  Let \(V\) and \(W\) be vector spaces over the same field \(\mathbb F\).
  A map \(L:V\to W\) is linear if it preserves addition and scalar
  multiplication:
  \[
    L(v+w)=L(v)+L(w),\qquad L(av)=aL(v)
  \]
  for \(v,w\in V\) and \(a\in\mathbb F\). Equivalently,
  \[
    L(av+bw)=aL(v)+bL(w)
  \]
  for all \(a,b\in\mathbb F\) and \(v,w\in V\).
  When \(W=\mathbb F\), a linear map is called a linear functional.
\end{definition}

In particular, \(L(0_V)=0_W\).
% September 24 slide 18 says linearity "implies" that the domain and
% codomain are vector spaces. Their vector-space structures are hypotheses
% here, as in the September 17 and quiz-prep definitions.

\begin{example}
  If \(E=(e_1,e_2)\) is a basis of \(V\), write
  \(a_1=L(e_1)\), \(a_2=L(e_2)\) for a linear functional \(L\).
  For \(v=c_1e_1+c_2e_2\),
  \[
    L(v)=c_1L(e_1)+c_2L(e_2)=c_1a_1+c_2a_2.
  \]
\end{example}

\subsection{Prescribing a map on a basis}

Let \((v_1,\ldots,v_m)\) be a basis of \(V\). Choose arbitrary vectors
\(w_1,\ldots,w_m\) in \(W\). The prescription
\[
  L(v_j)=w_j,\qquad 1\leq j\leq m,
\]
determines a linear map on all of \(V\). For
\(v=a^1v_1+\cdots+a^mv_m\), its value is
\[
  L(v)=a^1w_1+\cdots+a^mw_m.
\]

\begin{example}
  Define \(L:\R^2\to\R^2\) by
  \[
    L(v^1,v^2)=(3v^1-2v^2,\;5v^1+6v^2).
  \]
  In matrix notation,
  \[
    L(v)=\begin{bmatrix}3&-2\\5&6\end{bmatrix}
    \begin{bmatrix}v^1\\v^2\end{bmatrix}.
  \]
  In particular, \(L(1,1)=(1,11)\) and \(L(-2,3)=(-12,8)\).
\end{example}

\begin{example}
  The map \(T:\R^2\to\R^3\) given by
  \[
    T(v)=\begin{bmatrix}2&-1\\1&1\\-1&3\end{bmatrix}v
  \]
  sends the standard basis vectors to
  \(T(e_1)=(2,1,-1)\) and \(T(e_2)=(-1,1,3)\).
\end{example}

\subsection{Differentiation}

Differentiation is an example of a linear map whose inputs and outputs
are functions. Its definition, linearity, and application to polynomial
spaces are developed in
Section~\ref{sec:linear-algebra-i:differentiation}.
% The author requested a dedicated Chapter 4 section on October 1, 2026.
% The supplied September 17 slide 21 is retained there, with this pointer
% preserving its place in the introduction to linear maps.

\subsection{The vector space of linear maps}

The linear maps \(V\to W\) form a vector space, denoted
\(\mathcal L(V,W)\), under pointwise operations:
\[
  (L_1+L_2)(v)=L_1(v)+L_2(v),\qquad (aL)(v)=aL(v).
\]
If \(\dim V=m\) and \(\dim W=n\), choosing bases identifies this space
with the space of \(n\)-by-\(m\) matrices. The correspondence is developed
in Section~\ref{sec:linear-algebra-i:change-of-basis-for-functions-and-maps}.
